\documentclass[11pt]{article}
\usepackage[margin=0.9in]{geometry}
\usepackage{amsmath,amssymb,amsthm,mathtools}
\usepackage[T1]{fontenc}
\usepackage[utf8]{inputenc}
\usepackage{enumitem}
\usepackage{booktabs,longtable,array,seqsplit}
\usepackage{hyperref}
\usepackage{xcolor}
\usepackage{microtype}
\hypersetup{colorlinks=true,linkcolor=blue!45!black,citecolor=blue!45!black,urlcolor=blue!45!black}
\setlength{\emergencystretch}{6em}

\newtheorem{theorem}{Theorem}[section]
\newtheorem{proposition}[theorem]{Proposition}
\newtheorem{lemma}[theorem]{Lemma}
\newtheorem{corollary}[theorem]{Corollary}
\newtheorem{claim}[theorem]{Claim}
\theoremstyle{definition}
\newtheorem{definition}[theorem]{Definition}
\newtheorem{assumption}[theorem]{Assumption}
\theoremstyle{remark}
\newtheorem{remark}[theorem]{Remark}

\newcommand{\R}{\mathbb{R}}
\newcommand{\C}{\mathbb{C}}
\newcommand{\N}{\mathbb{N}}
\newcommand{\Hsp}{\mathcal{H}}
\newcommand{\Fsp}{\mathcal{F}}
\newcommand{\Id}{\mathrm{Id}}
\newcommand{\tr}{\operatorname{tr}}
\newcommand{\diag}{\operatorname{diag}}
\newcommand{\rank}{\operatorname{rank}}
\newcommand{\spec}{\operatorname{spec}}
\newcommand{\spanop}{\operatorname{span}}
\newcommand{\dist}{\operatorname{dist}}
\newcommand{\range}{\operatorname{ran}}
\newcommand{\Ker}{\operatorname{ker}}
\newcommand{\Pinch}{\mathcal{C}}
\newcommand{\Off}{\widetilde{\mathcal{C}}}
\newcommand{\Pperp}{P^{\perp}}
\newcommand{\Qperp}{Q^{\perp}}
\newcommand{\op}{\mathrm{op}}
\newcommand{\F}{\mathrm{F}}
\newcommand{\HS}{\mathrm{HS}}
\newcommand{\ip}[2]{\left\langle #1,#2\right\rangle}
\newcommand{\norm}[1]{\left\lVert #1\right\rVert}
\newcommand{\abs}[1]{\left\lvert #1\right\rvert}
\newcommand{\code}[1]{\nolinkurl{#1}}

\newenvironment{argument}{\paragraph{Argument.}\begin{enumerate}[leftmargin=2em]}{\end{enumerate}}
\newenvironment{proofoutline}{\paragraph{Proof architecture.}\begin{enumerate}[leftmargin=2em]}{\end{enumerate}}
\newcommand{\sourceanchor}[1]{\par\smallskip\noindent\textbf{Source anchor.} #1\par\smallskip}
\newcommand{\sourcecomment}[1]{\begin{remark}#1\end{remark}}
\newenvironment{sourceguide}{\begin{description}[style=nextline,leftmargin=3.3cm,labelwidth=3.1cm]}{\end{description}}
\newenvironment{formalizationmap}{\begin{longtable}{@{}p{0.20\textwidth}p{0.31\textwidth}p{0.40\textwidth}@{}}\toprule Source item & Modern mathematical content & Repository status \\\midrule\endfirsthead\toprule Source item & Modern mathematical content & Repository status \\\midrule\endhead}{\bottomrule\end{longtable}}

\title{The $\tan\Theta$ theorem after Davis--Kahan:\\graph subspaces, Riccati equations, relaxed conditions, and corrected finite statements}
\author{}
\date{Comparative formalization dossier}
\begin{document}
\maketitle

\section*{Sources and purpose}
\begin{sourceguide}
\item[Primary sources] V. Kostrykin, K. A. Makarov, and A. K. Motovilov, \emph{On the Existence of Solutions to the Operator Riccati Equation and the $\tan\Theta$ Theorem}, Integral Equations Operator Theory 51 (2005), 121--140, arXiv:math/0210032. Yuji Nakatsukasa, \emph{The $\tan\theta$ Theorem with Relaxed Conditions}, Linear Algebra Appl. 436 (2012), 1521--1534. A. K. Motovilov, \emph{Comment on ``The $\tan\theta$ Theorem with Relaxed Conditions''}, arXiv:1204.4441.
\item[Companion source] S. Albeverio and A. K. Motovilov, \emph{The a priori $\tan\Theta$ Theorem for Spectral Subspaces}, Integral Equations Operator Theory 73 (2012), 413--430, arXiv:1012.1569.
\item[Purpose] Separate four statements often conflated under the name ``tan theta theorem'': a posteriori residual bounds, graph/Riccati bounds, a priori spectral-subspace bounds, and relaxed Ritz-subspace conditions.
\item[Formalization target] \code{DavisKahan/Experimental/FiniteDimensional/TanTheta/GraphOperator.lean}, \code{DavisKahan/Experimental/InfiniteDimensional/GraphSubspace.lean}, \code{Riccati.lean}, and the generalized trial-map layer.
\end{sourceguide}

\section{Graph geometry is the primitive object}

Let $E=E_0\oplus E_1$ and let $X:E_0\to E_1$ be bounded. Its graph is
\[
  \mathcal G(X)=\{u\oplus Xu:u\in E_0\}.
\]
The graph is transverse to $E_1$ and the coordinate projection onto $E_0$ is invertible. The operator angle between $E_0$ and $\mathcal G(X)$ satisfies
\[
  \Theta=\arctan\sqrt{X^*X},
\]
so
\[
  \|\tan\Theta\|=\|X\|.
\]
In finite dimension the full singular-value statement is
\[
  s_j(X)=\tan\theta_j(E_0,\mathcal G(X)).
\]
This identity, rather than a totalized trigonometric operator, is the correct root for every-UI-norm tangent theorems.

\section{Block operator and Riccati equation}

Consider the self-adjoint block matrix
\[
  L=\begin{bmatrix}A&B\\B^*&C\end{bmatrix}
\]
on $E_0\oplus E_1$. Direct expansion shows that $\mathcal G(X)$ is invariant precisely when
\[
  XA-CX+XBX=B^*.
\]
For a self-adjoint $L$, an invariant closed graph is reducing. This yields the fundamental equivalence:
\[
  X\text{ solves Riccati}
  \quad\Longleftrightarrow\quad
  \mathcal G(X)\text{ reduces }L.
\]

Kostrykin--Makarov--Motovilov further use the graph transform
\[
  V=\begin{bmatrix}I&-X^*\\X&I\end{bmatrix}
\]
to obtain block diagonalization
\[
  V^{-1}LV=
  \begin{bmatrix}
    A+BX&0\\0&C-B^*X^*
  \end{bmatrix}.
\]
After normalization by $(I+X^*X)^{1/2}$ and $(I+XX^*)^{1/2}$, the diagonal blocks are similar to self-adjoint operators. This similarity is why their spectra remain real and why a Sylvester estimate can be applied to the Riccati equation rewritten around the perturbed diagonal block.

\section{Kostrykin--Makarov--Motovilov's generalized $\tan\Theta$ theorem}

Let $Z=A+BX$ be the diagonal block associated with the graph. Rewrite Riccati as
\[
  XZ-CX=B^*.
\]
Assume the spectrum of $Z$, equivalently the spectrum of $L|_{\mathcal G(X)}$, lies in a spectral gap of $C$, and let
\[
  \delta:=\dist(\spec(Z),\spec(C))>0.
\]
Then their Theorem 2 gives
\[
  \|X\|\le\frac{\|B\|}{\delta},
\]
or equivalently
\[
  \|\tan\Theta(E_0,\mathcal G(X))\|
  \le\frac{\|B\|}{\delta}.
\]

The important hypothesis is not merely $\dist(\spec(A),\spec(C))>0$. The denominator uses the spectrum of the \emph{perturbed graph block} $Z$. Thus the theorem is an a posteriori tangent theorem unless additional arguments control $\spec(Z)$ from unperturbed data.

\subsection{Proof mechanism}

Choose the center $\gamma$ of the spectral interval containing $\spec(Z)$. The gap hypothesis places $\spec(C)$ outside that interval by $\delta$. Rewriting
\[
  X(Z-\gamma I)-(C-\gamma I)X=B^*
\]
and multiplying by $(C-\gamma I)^{-1}$ gives
\[
  X=(C-\gamma I)^{-1}(X(Z-\gamma I)-B^*).
\]
The normalized self-adjoint representative of $Z$ yields a bound on $\|X(Z-\gamma I)\|$ with the same factor controlling $\|(C-\gamma I)^{-1}\|$. The scalar terms cancel to leave $\delta\|X\|\le\|B\|$.

In finite dimensions, an ordered rectangular Sylvester theorem gives a cleaner proof and automatically extends the estimate from operator norm to all rectangular UI norms, provided the graph equation has the correct left and right self-adjoint coefficients.

\section{Existence versus estimation}

The same paper proves existence and uniqueness of bounded Riccati solutions under stronger hypotheses. If $C$ has a finite spectral gap $\Delta$ containing $\spec(A)$ and
\[
  \|B\|<\sqrt{d|\Delta|},
  \qquad
  d=\dist(\spec(A),\spec(C)),
\]
then the spectral subspace of $L$ associated with $\Delta$ is the graph of a unique bounded solution with the prescribed spectral split.

Under additional boundedness and a stronger smallness condition, the solution is a strict contraction. The paper establishes sharp threshold phenomena: $\sqrt2$ is the optimal universal scale for existence in the finite-gap configuration, while $1$ is the sharp scale guaranteeing a contractive branch in the bounded setting.

These existence theorems should not be folded into a residual $\tan\Theta$ theorem. In finite DK work, the exact subspace may already be supplied and transversality may be derivable from the residual equation without a global Riccati existence theorem.

\section{Nakatsukasa's residual theorem as reconstructed by Motovilov}

Let $A$ be Hermitian. Let $Q_1$ have orthonormal columns and define its Ritz matrix and residual by
\[
  A_1=Q_1^*AQ_1,
  \qquad
  R=AQ_1-Q_1A_1.
\]
Let $X_1$ span an exact spectral subspace and let $\Lambda_2$ represent its complementary exact block. Suppose there is an interval $[\alpha,\beta]$ and $\delta>0$ such that
\[
  \spec(A_1)\subset(-\infty,\alpha-\delta]\cup[\beta+\delta,\infty),
\]
\[
  \spec(\Lambda_2)\subset[\alpha,\beta].
\]
Then Nakatsukasa's Theorem 1 states, in spectral norm,
\[
  \tan\angle(\range Q_1,\range X_1)
  \le\frac{\|R\|}{\delta}.
\]

Motovilov observes that this is a coordinate reformulation of the earlier graph theorem. Complete $Q_1$ to a unitary $Q=[Q_1\ Q_2]$ and set $L=Q^*AQ$. Then
\[
  L=\begin{bmatrix}A_1&B^*\\B&A_2\end{bmatrix},
  \qquad B=Q_2^*R,
\]
so $\|B\|=\|R\|$. The exact complementary spectral subspace lies in the gap of $A_1$. A dimension and separation argument proves the needed graph transversality automatically. Applying the earlier block $\tan\Theta$ theorem gives the result.

\section{What ``relaxed conditions'' actually change}

The relaxed result does not remove all separation. It changes which blocks are required to be separated.

Classical presentations often demand separation between the desired exact eigenvalue cluster and all approximate Ritz values. Nakatsukasa's formulation uses the approximate Ritz block $A_1$ and the \emph{unwanted exact complementary block} $\Lambda_2$. This is useful because the residual equation naturally couples those two blocks.

The theorem does not require prior knowledge that the approximate and exact subspaces are acute. Motovilov's Lemma 3 proves graph transversality from:
\begin{enumerate}
\item equal dimensions;
\item confinement of the exact complementary spectrum to $[\alpha,\beta]$;
\item exclusion of the approximate Ritz spectrum from the enlarged interval.
\end{enumerate}
A vector in the forbidden intersection would simultaneously satisfy incompatible lower and upper bounds for $\|(L-cI)y\|$, where $c=(\alpha+\beta)/2$.

This is a highly valuable Lean seam:
\[
  \text{spectral gap plus reducing exact block}
  \Longrightarrow
  \text{transversality}.
\]
It should be proved before the UI-norm tangent inequality.

\section{Motovilov's correction and clarification}

Motovilov's comment makes two points.

\begin{enumerate}
\item In spectral norm, Nakatsukasa's Theorem 1 is not a new independent analytic estimate; it is a finite-dimensional corollary of the 2005 graph/Riccati tangent theorem.
\item A second a priori tangent statement can be formulated using a residual threshold $\|R\|<\sqrt2d$, a spectral continuation result, and the same sharp bound
\[
  \tan\angle(Q_1,X_1)\le\frac{\|R\|}{d}.
\]
The threshold cannot in general be omitted because it is what guarantees existence and identification of the desired spectral branch.
\end{enumerate}

The correction is conceptual rather than a claim that the displayed residual inequality is false. It separates an a posteriori theorem, where the target exact spectral block is already identified, from an a priori theorem, where the existence and selection of that block require a smallness condition.

\section{Every-UI-norm finite parent theorem}

The literature statements above are mostly in operator norm. The maximal finite theorem should instead use the graph singular-value dictionary and rectangular Sylvester estimates.

Let $U$ reduce a self-adjoint $A$, let $X:F\to E$ be isometric, and let $M=M^*$. Assume:
\begin{enumerate}
\item the residual is $R=AX-XM$;
\item the spectral placement implies transversality of $\range X$ relative to $U^\perp$;
\item the graph map $T:F\to U^\perp$ satisfies an ordered Sylvester equation with gap $\delta$.
\end{enumerate}
Then for every rectangular UI norm $N$,
\[
  \delta N(T)\le N(R).
\]
The proof should be:
\begin{enumerate}
\item derive transversality;
\item construct $T$ with the proof of invertibility;
\item derive the exact graph Sylvester equation;
\item apply the ordered rectangular theorem;
\item identify the singular values of $T$ with principal tangents.
\end{enumerate}

This theorem subsumes the spectral-norm residual statement, Frobenius and Ky Fan forms, and perturbation theorems obtained by taking $X$ to embed a reducing subspace of a perturbed operator.

\section{Nonorthonormal trial maps}

For an injective trial map $Y$ with Gram operator $G=Y^*Y$, set
\[
  X=YG^{-1/2}.
\]
The residual transforms by right multiplication with $G^{-1/2}$. Therefore the generalized tangent theorem should not need a new Riccati proof. Its only new ingredient is a lower frame bound
\[
  \varepsilon^2I\le G,
\]
which yields
\[
  N(RG^{-1/2})\le \varepsilon^{-1}N(R).
\]
The graph geometry is carried by $\range Y=\range X$.

\section{Formalization decisions}

\begin{enumerate}
\item Define a proof-carrying \code{graphOperatorOfTransverse}; retain totalized maps only as wrappers.
\item Prove \code{singularValues_graphOperator} from principal planes.
\item Prove transversality from the residual-gap configuration before the norm theorem.
\item State the a posteriori theorem separately from branch-existence theorems.
\item Keep ordered-gap constant one separate from arbitrary spectral separation.
\item Avoid using a scalar maximal-angle theorem as the parent of all UI norms.
\item Make the Ritz/Galerkin hypothesis explicit when $M=X^*AX$ is used.
\end{enumerate}

\section{Repository map}

\begin{formalizationmap}
Graph angle identity & $\tan\Theta=|X|$ for graph subspaces & Experimental \code{GraphSubspace.lean}; key lemmas remain sorry. \\
Riccati equivalence & Reducing graph iff Riccati equation & Experimental \code{Riccati.lean}; block operator construction and equivalence unfinished. \\
A posteriori tangent & $\delta\|X\|\le\|B\|$ using perturbed graph spectrum & Operator-norm skeleton exists in older modules; maximal every-UI parent remains in \code{TanTheta.lean}. \\
Relaxed residual theorem & Ritz block separated from unwanted exact block & Signature exists; transversality and graph Sylvester equation are the hard proof seams. \\
A priori branch theorem & Residual smallness selects exact spectral branch & Depends on continuation and spectral inclusion modules. \\
Generalized trial map & Whitening transport & \code{Generalized.lean}; Gram inverse-square-root API remains incomplete. \\
\end{formalizationmap}

\section*{Bibliography}
\begin{thebibliography}{9}
\bibitem{KMM2005} V. Kostrykin, K. A. Makarov, and A. K. Motovilov, \emph{On the Existence of Solutions to the Operator Riccati Equation and the $\tan\Theta$ Theorem}, Integral Equations Operator Theory 51 (2005), 121--140.
\bibitem{Nakatsukasa2012} Y. Nakatsukasa, \emph{The $\tan\theta$ Theorem with Relaxed Conditions}, Linear Algebra Appl. 436 (2012), 1521--1534.
\bibitem{Motovilov2012} A. K. Motovilov, \emph{Comment on ``The $\tan\theta$ Theorem with Relaxed Conditions''}, arXiv:1204.4441.
\bibitem{AM2012} S. Albeverio and A. K. Motovilov, \emph{The a priori $\tan\Theta$ Theorem for Spectral Subspaces}, Integral Equations Operator Theory 73 (2012), 413--430.
\end{thebibliography}

\end{document}
